% ============================================================
%  sections/06_conclusion.tex
%  6. Kết luận & Hướng phát triển  (~0.5 trang)
% ============================================================

\section{Kết Luận và Hướng Phát Triển}
\label{sec:conclusion}

Bài báo này trình bày nghiên cứu so sánh hệ thống bốn phương pháp phát
hiện văn bản do máy tạo ra (MGT) trên dữ liệu mạng xã hội đa ngôn ngữ
gồm bốn ngôn ngữ: EN, VI, ZH, AR. Bốn kết luận chính được rút ra:

\begin{enumerate}
    \item \textbf{Thiên lệch Anh ngữ là nghiêm trọng và nguy hiểm:}
    Bộ phát hiện tiền huấn luyện chỉ dành cho tiếng Anh (\texttt{roberta-base-openai-detector})
    đạt dưới mức ngẫu nhiên trên VI, ZH, AR --- thậm chí phân loại sai
    hệ thống với AR (F1 = 0,336). Triển khai detector đơn ngôn ngữ cho
    bài toán đa ngôn ngữ là không thể chấp nhận được.

    \item \textbf{Tinh chỉnh XLM-RoBERTa là giải pháp thực tế tối ưu:}
    Đạt \textbf{81,4\% accuracy} trung bình, huấn luyện trong 7,5
    phút trên GPU T4 thông thường với chi phí ước tính \textbf{\$1,50}.
    Phương pháp này chứng minh rằng phát hiện MGT đa ngôn ngữ chất lượng
    cao không đòi hỏi phần cứng A100 đắt tiền.

    \item \textbf{PPL zero-shot hữu ích có điều kiện:}
    Phát hiện dựa trên perplexity (Qwen2.5-1.5B) hoạt động tốt khi tỉ
    lệ PPL giữa hai lớp đủ lớn (EN: 6,4$\times$, AR: 3,4$\times$), nhưng
    thất bại với văn bản MXH nhiễu (VI Teencode, ZH Weibo ngắn).

    \item \textbf{Học chuyển giao liên ngôn ngữ là chưa đủ:}
    Thí nghiệm leave-one-language-out cho thấy mỗi ngôn ngữ cần dữ liệu
    huấn luyện riêng. Tiếng Việt đặc biệt khó khăn (F1 = 0,134 khi
    không có dữ liệu VI trong tập huấn luyện), khẳng định tính tất yếu
    của dữ liệu đa ngôn ngữ trong bài toán MGT trên MXH.
\end{enumerate}

\subsection*{Hướng Phát Triển}

Nghiên cứu tiếp theo có thể khai thác các hướng sau:

\begin{itemize}
    \item \textbf{Mô hình cấp ký tự} cho VI và ZH nhằm xử lý đặc tính
          Teencode và bài đăng siêu ngắn (15--30 ký tự).
    \item \textbf{Data augmentation} với Teencode giả lập để cải thiện
          khả năng tổng quát hóa trên VI.
    \item \textbf{Mô hình lớn hơn với quantization:} thử nghiệm XLM-R
          large (559M) hoặc mGPT với INT8/INT4 quantization trên T4.
    \item \textbf{Mở rộng ngôn ngữ:} Thêm tiếng Nhật, Hàn, Hindi ---
          các ngôn ngữ MXH lớn chưa được nghiên cứu kỹ trong bài toán MGT.
    \item \textbf{Phát hiện cấp đặc trưng ngôn ngữ học:} kết hợp PPL
          với đặc trưng cú pháp và ngữ nghĩa để vượt qua giới hạn của
          phương pháp thuần PPL.
\end{itemize}

Mã nguồn, bộ dữ liệu và pipeline thực nghiệm sẽ được công khai để
thúc đẩy nghiên cứu tái lập và mở rộng.
